% GENERATED FILE. Edit canonical lessons, then run npm run generate:books.
% canonical-source-hash: fnv1a64:05ac360d40ba4d19
% canonical-lessons: TE-C204-polcu, TE-C204-ralu, TE-C204-vyayamam-ceyu, TE-C204-natincu, TE-C204-giyu

\chapter{To compare, To fall, To exercise, To act, To draw}
\label{ch:te-to-compare-to-fall-to-exercise-to-act-to-draw}
\hlchaptermodality{204}

\begin{chapteropening}
\textbf{By the end of this chapter:} \emph{I can say to compare, to fall, to exercise, to act and to draw.}

Complete the last lesson of chapter 204: I can say to compare, to fall, to exercise, to act and to draw.
\end{chapteropening}

\section[\texorpdfstring{\te{పోల్చు} (pōlcu)}{పోల్చు (pōlcu)}]{\te{పోల్చు} (pōlcu) — to compare}

\label{lesson:TE-C204-polcu}



\begin{quote}
Before the new one: say the Telugu for to bloom, then the Telugu for to rain down.
\end{quote}

\subsection*{You'll want to know: \te{పోల్చు}}
\textbf{\te{పోల్చు}} — \emph{pōlcu} — \textquotedblleft{}to compare\textquotedblright{}.

\begin{grammarlens}[title={What you've built}]
One more word for this chapter.
\end{grammarlens}

\subsection*{Your turn}
\begin{itemize}
\raggedright
  \item \emph{Say it:} \emph{pōlcu}
  \item \emph{Say it:} \emph{pōlcu}, once more
  \item \emph{Say it:} say \emph{pōlcu}
  \item \emph{Recall:} say the Telugu for to bloom, then the Telugu for to rain down, then say \emph{pōlcu} again
\end{itemize}

\begin{tcolorbox}[breakable,colback=teal!4,colframe=teal!35!black,title={Before you move on}]
What does \te{పోల్చు} mean? (\textquotedblleft{}To compare\textquotedblright{}.) Say it once more.
\end{tcolorbox}
\section[\texorpdfstring{\te{రాలు} (rālu)}{రాలు (rālu)}]{\te{రాలు} (rālu) — to fall (of leaves, fruit, petals)}

\label{lesson:TE-C204-ralu}



\begin{quote}
Before the new one: say the Telugu for to rain down, then the Telugu for to compare.
\end{quote}

\subsection*{You'll want to know: \te{రాలు}}
\textbf{\te{రాలు}} — \emph{rālu} — \textquotedblleft{}to fall (of leaves, fruit, petals)\textquotedblright{}.

\begin{grammarlens}[title={What you've built}]
One more word for this chapter.
\end{grammarlens}

\subsection*{Your turn}
\begin{itemize}
\raggedright
  \item \emph{Say it:} \emph{rālu}
  \item \emph{Say it:} \emph{rālu}, once more
  \item \emph{Say it:} say \emph{rālu}
  \item \emph{Recall:} say the Telugu for to rain down, then the Telugu for to compare, then say \emph{rālu} again
\end{itemize}

\begin{tcolorbox}[breakable,colback=teal!4,colframe=teal!35!black,title={Before you move on}]
What does \te{రాలు} mean? (\textquotedblleft{}To fall (of leaves, fruit, petals)\textquotedblright{}.) Say it once more.
\end{tcolorbox}
\section[\texorpdfstring{\te{వ్యాయామం} \te{చేయు} (vyāyāmaṁ cēyu)}{వ్యాయామం చేయు (vyāyāmaṁ cēyu)}]{\te{వ్యాయామం} \te{చేయు} (vyāyāmaṁ cēyu) — to exercise}

\label{lesson:TE-C204-vyayamam-ceyu}



\begin{quote}
Before the new one: say the Telugu for to compare, then the Telugu for to fall.
\end{quote}

\subsection*{You'll want to know: \te{వ్యాయామం} \te{చేయు}}
\textbf{\te{వ్యాయామం} \te{చేయు}} — \emph{vyāyāmaṁ cēyu} — \textquotedblleft{}to exercise\textquotedblright{}.

\begin{grammarlens}[title={What you've built}]
One more word for this chapter.
\end{grammarlens}

\subsection*{Your turn}
\begin{itemize}
\raggedright
  \item \emph{Say it:} \emph{vyāyāmaṁ cēyu}
  \item \emph{Say it:} \emph{vyāyāmaṁ cēyu}, once more
  \item \emph{Say it:} say \emph{vyāyāmaṁ cēyu}
  \item \emph{Recall:} say the Telugu for to compare, then the Telugu for to fall, then say \emph{vyāyāmaṁ cēyu} again
\end{itemize}

\begin{tcolorbox}[breakable,colback=teal!4,colframe=teal!35!black,title={Before you move on}]
What does \te{వ్యాయామం} \te{చేయు} mean? (\textquotedblleft{}To exercise\textquotedblright{}.) Say it once more.
\end{tcolorbox}
\section[\texorpdfstring{\te{నటించు} (naṭiñcu)}{నటించు (naṭiñcu)}]{\te{నటించు} (naṭiñcu) — to act (in a play or film)}

\label{lesson:TE-C204-natincu}



\begin{quote}
Before the new one: say the Telugu for to fall, then the Telugu for to exercise.
\end{quote}

\subsection*{You'll want to know: \te{నటించు}}
\textbf{\te{నటించు}} — \emph{naṭiñcu} — \textquotedblleft{}to act (in a play or film)\textquotedblright{}.

\begin{grammarlens}[title={What you've built}]
One more word for this chapter.
\end{grammarlens}

\subsection*{Your turn}
\begin{itemize}
\raggedright
  \item \emph{Say it:} \emph{naṭiñcu}
  \item \emph{Say it:} \emph{naṭiñcu}, once more
  \item \emph{Say it:} say \emph{naṭiñcu}
  \item \emph{Recall:} say the Telugu for to fall, then the Telugu for to exercise, then say \emph{naṭiñcu} again
\end{itemize}

\begin{tcolorbox}[breakable,colback=teal!4,colframe=teal!35!black,title={Before you move on}]
What does \te{నటించు} mean? (\textquotedblleft{}To act (in a play or film)\textquotedblright{}.) Say it once more.
\end{tcolorbox}
\section[\texorpdfstring{\te{గీయు} (gīyu)}{గీయు (gīyu)}]{\te{గీయు} (gīyu) — to draw (a line, a picture)}

\label{lesson:TE-C204-giyu}



\begin{quote}
Before the new one: say the Telugu for to exercise, then the Telugu for to act.
\end{quote}

\subsection*{You'll want to know: \te{గీయు}}
\textbf{\te{గీయు}} — \emph{gīyu} — \textquotedblleft{}to draw (a line, a picture)\textquotedblright{}.

\begin{grammarlens}[title={What you've built}]
One more word for this chapter.
\end{grammarlens}

\subsection*{Your turn}
\begin{itemize}
\raggedright
  \item \emph{Say it:} \emph{gīyu}
  \item \emph{Say it:} \emph{gīyu}, once more
  \item \emph{Say it:} say \emph{gīyu}
  \item \emph{Recall:} say the Telugu for to exercise, then the Telugu for to act, then say \emph{gīyu} again
\end{itemize}

\begin{tcolorbox}[breakable,colback=teal!4,colframe=teal!35!black,title={Before you move on}]
What does \te{గీయు} mean? (\textquotedblleft{}To draw (a line, a picture)\textquotedblright{}.) Say it once more.
\end{tcolorbox}
